\documentclass[a4paper,12pt]{article}
\usepackage[T2A]{fontenc}
\usepackage[utf8]{inputenc}
\usepackage[russian]{babel}
\usepackage{amsmath,amssymb,mathtools}
\usepackage{helvet}
\renewcommand{\familydefault}{\sfdefault}
\usepackage{icomma}
\usepackage[left=19mm,right=19mm,top=18mm,bottom=19mm]{geometry}
\usepackage{microtype}
\usepackage{tikz}
\usetikzlibrary{calc,arrows.meta,positioning,shapes.geometric}
\usepackage{tkz-euclide}
\usepackage{pgfplots}
\pgfplotsset{compat=1.18}
\usepackage{enumitem,tabularx,booktabs}
\usepackage{fancyhdr,parskip,setspace,needspace}
\usepackage{xcolor,hyperref}
\definecolor{accent}{HTML}{166D72}
\definecolor{darkaccent}{HTML}{15484F}
\definecolor{textgray}{HTML}{31383A}
\definecolor{softgray}{HTML}{667579}
\definecolor{amber}{HTML}{AC6A38}
\hypersetup{colorlinks=true,linkcolor=darkaccent,urlcolor=darkaccent}
\setlength{\parindent}{0pt}
\setlength{\parskip}{4pt}
\setstretch{1.2}
\setlist[itemize]{leftmargin=7mm,itemsep=1.5pt,topsep=2pt}
\setlist[enumerate]{leftmargin=9mm,itemsep=5pt,topsep=3pt}
\pagestyle{fancy}
\fancyhf{}
\fancyfoot[C]{\small\color{softgray}\thepage}
\renewcommand{\headrulewidth}{0pt}
\newcommand{\sect}[1]{\par\vspace{2.5mm}{\Large\bfseries\color{darkaccent}#1}\par\vspace{1mm}{\color{accent!50}\rule{\linewidth}{.55pt}}\par\vspace{1mm}}
\newcommand{\key}[1]{\textbf{\color{darkaccent}#1}}
\newcommand{\ck}[1]{\par\noindent$\square$\enspace #1\par}
\newcommand{\ex}[1]{\par\smallskip\noindent\textbf{#1.}\enspace}
\begin{document}
\thispagestyle{empty}
\begin{center}
\vspace*{31mm}
{\small\color{softgray}МАТЕМАТИКА \textbullet{} ОГЭ \textbullet{} ПОВТОРЕНИЕ И ЗАДАЧИ ПОВЫШЕННОЙ СЛОЖНОСТИ}\par
\vspace{14mm}
{\fontsize{30}{34}\selectfont\bfseries\color{darkaccent}Чек-лист}\par
\vspace{7mm}
{\LARGE\bfseries Уравнения, движение, графики и подобие}\par
\vspace{4mm}
{\large Четыре метода решения задач}\par
\vspace{19mm}
{\large София Кальней}\par
\vspace{4mm}
{\normalsize 18 октября 2026 года}\par
\vfill
{\small\color{textgray}Лёвин Артём Александрович эксклюзивно для Софии Кальней}\par
\vspace{17mm}
\end{center}
\begin{tikzpicture}[remember picture,overlay]
\draw[accent!50,line width=1.2pt]
 ([xshift=9mm,yshift=8mm]current page.south west)
 .. controls +(45mm,23mm) and +(-54mm,-6mm) ..
 ([xshift=-9mm,yshift=16mm]current page.south east);
\draw[amber!45,line width=.65pt]
 ([xshift=15mm,yshift=6mm]current page.south west)
 .. controls +(52mm,18mm) and +(-65mm,-8mm) ..
 ([xshift=-15mm,yshift=12mm]current page.south east);
\end{tikzpicture}
\clearpage

\sect{1. Уравнение: разность кубов и неотрицательность}
\ck{Узнаю формулу $(a^3-b^3)=(a-b)(a^2+ab+b^2)$ и правильно определяю $a$ и $b$.}
\ck{Использую неотрицательность чётной степени и сохранение знака при возведении в куб.}
\ck{Если сумма неотрицательных выражений равна нулю, записываю \emph{систему} условий равенства нулю.}

\key{Пример с занятия.} Решим $x^6=(8-7x)^3$. Положим $a=x^2$, $b=8-7x$. Тогда
\[
 (a-b)(a^2+ab+b^2)=0.
\]
Второй множитель всегда неотрицателен, поскольку
\[
 a^2+ab+b^2=\left(a+\frac b2\right)^2+\frac{3b^2}{4}.
\]
Он равен нулю \emph{только} при $a=b=0$. Здесь это потребовало бы одновременно $x^2=0$ и $8-7x=0$, что невозможно. Следовательно, остаётся
\[
 x^2-(8-7x)=0\quad\Longleftrightarrow\quad
 x^2+7x-8=(x+8)(x-1)=0.
\]
\key{Ответ:} $x=-8$ или $x=1$.

\key{Почему работает рассуждение со знаком.} Из $x^6\ge0$ и исходного равенства следует $(8-7x)^3\ge0$, то есть $8-7x\ge0$ \emph{для искомых решений}. Поэтому три слагаемых $x^4$, $x^2(8-7x)$, $(8-7x)^2$ неотрицательны. Равенство их суммы нулю потребовало бы системы
\[
\begin{cases}
 x^4=0,\\x^2(8-7x)=0,\\(8-7x)^2=0,
\end{cases}
\]
которая несовместна. В этом рассуждении важно указывать, что все условия должны выполняться \emph{одновременно}.

\key{Короткий путь.} Для действительных чисел $a^3=b^3\iff a=b$, потому что функция $t^3$ строго возрастает. Поэтому можно сразу перейти от $x^6=(8-7x)^3$ к $x^2=8-7x$. Оба способа дают один ответ.

\clearpage
\sect{2. Движение по реке: таблица и отсрочка старта}
\ck{Отличаю собственную скорость лодки $x$ от скорости течения $u$; обязательно требую $x>u$.}
\ck{Записываю скорости: по течению $x+u$, против течения $x-u$, плот --- $u$.}
\ck{Заполняю таблицу $s,v,t$ по формуле $t=s/v$; сравниваю время с учётом момента старта.}

\key{Пример с занятия.} Скорость течения $3$ км/ч. Лодка проходит по $108$ км по течению и против него. Плот вышел на $2$ ч раньше лодки и к моменту её возвращения прошёл $51$ км. Найдём собственную скорость лодки $x$.
\begin{center}
\renewcommand{\arraystretch}{1.52}
\begin{tabularx}{\linewidth}{@{}Xccc@{}}
\toprule
Участник / участок & $s$, км & $v$, км/ч & $t$, ч\\
\midrule
Лодка по течению & $108$ & $x+3$ & $\tfrac{108}{x+3}$\\[1mm]
Лодка против течения & $108$ & $x-3$ & $\tfrac{108}{x-3}$\\[1mm]
Плот & $51$ & $3$ & $17$\\
\bottomrule
\end{tabularx}
\end{center}
Плот был в пути $17$ ч, лодка --- $17-2=15$ ч. Поэтому
\[
\frac{108}{x+3}+\frac{108}{x-3}=15,\quad x>3.
\]
После приведения к общему знаменателю получаем
\[
 5x^2-72x-45=0,\qquad x=15\ \text{или}\ x=-\frac35.
\]
С учётом $x>3$ ответ: \key{$15$ км/ч}.

\key{Проверка найденного ответа.} При $x=15$ лодка затрачивает $108/18=6$ ч по течению и $108/12=9$ ч против течения. Всего $15$ ч. Плот за $17$ ч проходит $51$ км; разница длительностей действительно $2$ ч.

\key{Общий шаблон.} Если течение имеет скорость $u$, плот прошёл $S$ км, вышел на $\Delta t$ часов раньше лодки и оба закончили движение одновременно, а лодка прошла по $L$ км в каждом направлении, то
\[
\boxed{\frac{L}{x+u}+\frac{L}{x-u}=\frac Su-\Delta t},\qquad x>u.
\]
Сначала проверяйте, что $S/u-\Delta t>0$, затем решайте полученное рациональное уравнение и отбрасывайте корни вне допустимой области.

\key{Типичные ошибки.}
\begin{itemize}
\item Приравнивают время лодки ко времени плота, забывая, что участники стартовали в разные моменты.
\end{itemize}
\clearpage

\sect{3. График с модулем и число пересечений}
\ck{Сначала нахожу нули подмодульного выражения и интервалы его знака.}
\ck{На каждом интервале раскрываю модуль со своим знаком; строю только допустимую часть полученной функции.}
\ck{Для прямой $y=m$ считаю \emph{различные точки} пересечения; значение $m$ --- ордината, а не абсцисса.}

\key{Пример с занятия.} Исследуем
\[
 f(x)=x^2-2x-|x^2-4|-3.
\]
Знак $x^2-4$ меняется в точках $-2$ и $2$, поэтому
\[
 f(x)=\begin{cases}
 -2x+1,&x\le-2\ \text{или}\ x\ge2,\\
 2x^2-2x-7,&-2<x<2.
 \end{cases}
\]
Точки стыка: $(-2;5)$ и $(2;-3)$. Вершина внутренней параболы:
\[
 x_0=-\frac{b}{2a}=\frac12,\qquad
 y_0=f\!\left(\frac12\right)=-\frac{15}{2}.
\]
\begin{center}
\begin{tikzpicture}
\begin{axis}[
 width=12.3cm,height=8.0cm,
 axis lines=middle,axis line style={softgray},
 xmin=-5.3,xmax=6.2,ymin=-13,ymax=12,
 xtick={-4,-2,0,2,4,6},ytick={-10,-7.5,-3,0,5,10},
 tick label style={font=\scriptsize,text=softgray},
 xlabel={$x$},ylabel={$y$},xlabel style={at={(axis description cs:1.01,0.46)},anchor=west},
 ylabel style={at={(axis description cs:.46,1.01)},anchor=south},
 grid=both,major grid style={draw=softgray!15},minor grid style={draw=softgray!8},minor tick num=1,
 clip=false
]
\addplot[domain=-5.3:-2,samples=80,very thick,color=accent] {-2*x+1};
\addplot[domain=-2:2,samples=100,very thick,color=accent] {2*x*x-2*x-7};
\addplot[domain=2:6.2,samples=80,very thick,color=accent] {-2*x+1};
\addplot[domain=-5.3:6.2,samples=2,dashed,thin,color=amber!85] {-7.5};
\addplot[domain=-5.3:6.2,samples=2,dashed,thin,color=amber!85] {-3};
\addplot[only marks,mark=*,mark size=2pt,color=darkaccent] coordinates {(-2,5)(2,-3)(0.5,-7.5)};
\node[font=\scriptsize,anchor=south west,color=darkaccent] at (axis cs:-2,5) {$(-2;5)$};
\node[font=\scriptsize,anchor=south west,color=darkaccent] at (axis cs:2,-3) {$(2;-3)$};
\node[font=\scriptsize,anchor=north west,color=darkaccent] at (axis cs:0.5,-7.5) {$\left(\frac12;-\frac{15}{2}\right)$};
\end{axis}
\end{tikzpicture}
\end{center}
\key{Когда ровно две общие точки?} При $m=-\frac{15}{2}$ горизонталь касается вершины параболы и пересекает правый луч. При $m=-3$ она пересекает параболу в $x=-1$ и проходит через точку стыка $x=2$. Других значений нет:
\[
\boxed{m\in\left\{-\frac{15}{2},\,-3\right\}.}
\]
\key{Контроль.} $x_0=\frac12$ --- \emph{абсцисса} вершины; значение параметра в горизонтальной прямой $y=m$ равно её \emph{ординате} $-\frac{15}{2}$. В точках стыка график непрерывен, они не выколоты.
\clearpage

\sect{4. Подобие треугольников в трапеции}
\ck{Использую параллельность оснований для нахождения равных накрест лежащих углов.}
\ck{Проверяю равенство \emph{двух отношений сторон}, заключающих эти углы.}
\ck{Записываю соответствие вершин и только после этого --- коэффициент подобия.}

\key{Пример с занятия.} В трапеции $ABCD$ основания $BC\parallel AD$, причём $BC=9$, $BD=18$, $AD=36$. Докажем, что $\triangle CBD\sim\triangle BDA$.
\begin{center}
\begin{tikzpicture}[scale=1.05]
\tkzDefPoint(0,0){A}
\tkzDefPoint(4.8,2.683282){B}
\tkzDefPoint(6.6,2.683282){C}
\tkzDefPoint(7.2,0){D}
\tkzDrawPolygon[draw=textgray,line width=1pt](A,B,C,D)
\tkzDrawSegment[draw=accent,line width=1.3pt](B,D)
\tkzDrawPoints[size=2.3,fill=darkaccent](A,B,C,D)
\tkzLabelPoint[below left](A){$A$}
\tkzLabelPoint[above left](B){$B$}
\tkzLabelPoint[above right](C){$C$}
\tkzLabelPoint[below right](D){$D$}
\tkzLabelSegment[above](B,C){$9$}
\tkzLabelSegment[below](A,D){$36$}
\tkzLabelSegment[below left](B,D){$18$}
\end{tikzpicture}
\end{center}
Из $BC\parallel AD$ следует $\angle CBD=\angle BDA$. При этом
\[
\frac{BC}{BD}=\frac9{18}=\frac12,
\qquad \frac{BD}{DA}=\frac{18}{36}=\frac12.
\]
Значит, стороны, заключающие равные углы, пропорциональны. По второму признаку подобия (две пропорциональные стороны и угол между ними)
\[
\boxed{\triangle CBD\sim\triangle BDA}.
\]
\key{Соответствие:} $C\leftrightarrow B$, $B\leftrightarrow D$, $D\leftrightarrow A$. Например, $\frac{CD}{BA}=\frac{BC}{BD}=\frac12$.

\sect{5. Перед самостоятельной работой}
\ck{При разности кубов проверяю, может ли второй множитель обращаться в нуль.}
\ck{В речных задачах пишу $x>u$ и отдельно отмечаю временную задержку.}
\ck{У модуля правильно решаю неравенство, включая или исключая граничные точки.}
\ck{Для $y=m$ считаю количество пересечений, а не количество ветвей или корней на глаз.}
\ck{В доказательстве подобия сначала даю основания (углы, отношения), затем вывод.}
\clearpage

\begin{center}
{\Large\bfseries\color{darkaccent}Алгоритм: построить график с модулем}\par
\vspace{3mm}
{\small\color{softgray}Пример $f(x)=x^2-2x-|x^2-4|-3$}
\end{center}
\vspace{8mm}
\begin{center}
\begin{tikzpicture}[
 >=Stealth,node distance=9mm,
 start/.style={ellipse,draw=darkaccent,line width=.8pt,minimum height=10mm,align=center,text width=7.5cm,fill=accent!4},
 flowstep/.style={rectangle,rounded corners=1.3mm,draw=accent!85,line width=.85pt,align=center,text width=10cm,minimum height=11mm,inner sep=2.6mm},
 choose/.style={diamond,draw=amber,line width=.85pt,aspect=2.6,align=center,text width=4.3cm,inner sep=1mm},
 branch/.style={rectangle,rounded corners=1mm,draw=accent!70,line width=.8pt,align=center,text width=5cm,minimum height=21mm,inner sep=2.5mm},
 arr/.style={->,draw=darkaccent,line width=1pt}
]
\node[start](a){Исходная функция с модулем};
\node[flowstep,below=of a](b){Приравнять подмодульное выражение к нулю:\\ $x^2-4=0\Rightarrow x=\pm2$};
\node[choose,below=11mm of b](d){$x^2-4\ge0$?};
\node[branch,below=18mm of d,xshift=-37mm](l){\textbf{Да:} $x\le-2$ или $x\ge2$\\ $|x^2-4|=x^2-4$\\ $f(x)=-2x+1$};
\node[branch,below=18mm of d,xshift=37mm](r){\textbf{Нет:} $-2<x<2$\\ $|x^2-4|=4-x^2$\\ $f(x)=2x^2-2x-7$};
\coordinate(mid) at ($(l.south)!0.5!(r.south)$);
\node[flowstep,below=16mm of mid](e){Построить \emph{только нужные части} прямой и параболы; отметить точки стыка и вершину};
\node[flowstep,below=of e](f){Провести $y=m$ и сосчитать различные точки пересечения на каждом характерном уровне};
\node[start,below=of f](g){Записать значения $m$};
\draw[arr](a)--(b);
\draw[arr](b)--(d);
\draw[arr](d.south)--++(0,-6mm)-|node[pos=.86,above,font=\scriptsize]{да}(l.north);
\draw[arr](d.south)--++(0,-6mm)-|node[pos=.86,above,font=\scriptsize]{нет}(r.north);
\coordinate(join) at ($(e.north)+(0,4mm)$);
\draw[draw=darkaccent,line width=1pt](l.south)|-(join);
\draw[draw=darkaccent,line width=1pt](r.south)|-(join);
\draw[arr](join)--(e.north);
\draw[arr](e)--(f);
\draw[arr](f)--(g);
\end{tikzpicture}
\end{center}
\clearpage

\sect{6. Самостоятельная тренировка --- Путь А}
\textbf{10 заданий.} Располагайте преобразования последовательно. В речных задачах составьте таблицу скоростей, расстояний и времени. В геометрии сделайте рисунок и обоснуйте подобие.

\textbf{Уравнения.}
\begin{enumerate}[leftmargin=9mm,itemsep=4mm]
\item Решите уравнение $x^6=(8-7x)^3$ двумя способами: через разность кубов и через равенство оснований кубов.
\item Решите уравнение $x^6=(1-2x)^3$. Объясните, почему решения второго множителя при разложении разности кубов отсутствуют.
\item Решите уравнение $x^6=(4-x)^3$. Запишите корни в точном виде.
\end{enumerate}

\textbf{Движение по реке.}
\begin{enumerate}[start=4,leftmargin=9mm,itemsep=4mm]
\item Собственная скорость моторной лодки неизвестна, скорость течения $3$ км/ч. Лодка прошла $45$ км по течению и $45$ км против течения за $8$ ч. Найдите собственную скорость лодки.
\item Скорость течения реки $2$ км/ч. Плот отправился на $4$ ч раньше лодки и к моменту её возвращения прошёл $28$ км. Лодка прошла $48$ км вниз и $48$ км вверх по реке. Найдите скорость лодки в стоячей воде.
\item Скорость течения $2$ км/ч. Плот начал движение на $3$ ч раньше лодки; к моменту её возвращения плот прошёл $34$ км. Лодка прошла $96$ км по течению и столько же против течения. Найдите её собственную скорость.
\end{enumerate}

\textbf{Графики и параметр.}
\begin{enumerate}[start=7,leftmargin=9mm,itemsep=4mm]
\item Постройте график $y=x^2-2x-|x^2-4|-3$. Найдите все значения $m$, при которых прямая $y=m$ имеет с графиком ровно две общие точки.
\item Постройте график $y=x^2-|x^2-4|$. Для каких значений $m$ прямая $y=m$ пересекает его ровно в двух точках? Отдельно исследуйте уровень $m=4$.
\end{enumerate}

\textbf{Геометрия.}
\begin{enumerate}[start=9,leftmargin=9mm,itemsep=4mm]
\item В трапеции $ABCD$ основания $BC\parallel AD$, $BC=9$, $BD=18$, $AD=36$. Докажите, что $\triangle CBD\sim\triangle BDA$. Укажите коэффициент подобия первого треугольника ко второму.
\item В трапеции $ABCD$ основания $BC\parallel AD$. Известно, что $BC=12$, $BD=18$, $AD=27$, $CD=8$. Докажите подобие $\triangle CBD$ и $\triangle BDA$, затем найдите сторону $AB$.
\end{enumerate}
\clearpage

\sect{7. Ответы и самопроверка}
\renewcommand{\arraystretch}{1.4}
\begin{tabularx}{\linewidth}{@{}cX@{}}
\toprule
№ & Ответ / ключевой результат\\
\midrule
1 & $x=-8$; $x=1$.\\
2 & $x=-1-\sqrt2$; $x=-1+\sqrt2$.\\
3 & $x=\dfrac{-1-\sqrt{17}}2$; $x=\dfrac{-1+\sqrt{17}}2$.\\
4 & $12$ км/ч, допустимая область $x>3$.\\
5 & $10$ км/ч. Плот двигался $14$ ч, лодка --- $10$ ч.\\
6 & $14$ км/ч. Плот двигался $17$ ч, лодка --- $14$ ч.\\
7 & $m=-\dfrac{15}{2}$ и $m=-3$.\\
8 & $-4<m<4$. При $m=4$ общих точек бесконечно много (оба внешних луча лежат на горизонтали $y=4$).\\
9 & $\triangle CBD\sim\triangle BDA$; коэффициент подобия первого треугольника ко второму --- $\dfrac12$.\\
10 & $AB=12$. По подобию $\dfrac{CD}{AB}=\dfrac{BC}{BD}=\dfrac23$.\\
\bottomrule
\end{tabularx}

\vspace{7mm}
\key{Контрольные соотношения для задач 4--6.}
\[
\begin{aligned}
4:\quad& \frac{45}{x+3}+\frac{45}{x-3}=8,\qquad x>3;\\[2mm]
5:\quad& \frac{48}{x+2}+\frac{48}{x-2}=\frac{28}{2}-4=10,\quad x>2;\\[2mm]
6:\quad& \frac{96}{x+2}+\frac{96}{x-2}=\frac{34}{2}-3=14,\quad x>2.
\end{aligned}
\]

\key{Контрольные свойства графика из задачи 7.} Вершина внутренней параболы имеет координаты $\left(\frac12;-\frac{15}{2}\right)$, а точки стыка --- $(-2;5)$ и $(2;-3)$. Эти сведения позволяют независимо проверить оба значения $m$.

\vfill
{\small\color{softgray}При выполнении геометрических заданий важна последовательность доказательства: сначала равные углы и отношения сторон, затем признак подобия и вывод.}
\end{document}
